\documentclass{article}
\usepackage{amsmath,amssymb,amsthm}
\usepackage{algorithm}
\usepackage{algpseudocode}
\usepackage{hyperref}
\usepackage{enumitem}
\usepackage{bm}
\usepackage{geometry}
\geometry{margin=1in}
\newtheorem{theorem}{Theorem}
\newtheorem{lemma}{Lemma}
\newtheorem{assumption}{Assumption}
\newtheorem{corollary}{Corollary}
\newcommand{\E}{\mathbb{E}}
\newcommand{\R}{\mathbb{R}}
\newcommand{\norm}[1]{\left\lVert #1\right\rVert}
\newcommand{\ip}[2]{\langle #1,#2\rangle}
\newcommand{\grad}{\nabla}
\begin{document}

\section*{Heavy–Ball Method (HB)}

\section*{Mathematical Formulation}
Let $f:\R^{d}\rightarrow\R$ be differentiable.  
Given an initial pair $(x_{0},x_{-1})\in\R^{d}\times\R^{d}$, step size $\alpha>0$ and momentum parameter $\beta\in[0,1)$, the Heavy–Ball iteration is
\begin{equation}
\boxed{
x_{k+1}=x_{k}+\beta\,(x_{k}-x_{k-1})-\alpha\,\grad f(x_{k}),\qquad k=0,1,2,\dots
}
\label{eq:HB}
\end{equation}
The term $\beta\,(x_{k}-x_{k-1})$ incorporates the previous displacement (momentum).  
When $f$ is $\mu$‑strongly convex and $L$‑smooth (see Assumptions), the method enjoys a linear convergence rate for a suitable choice of $(\alpha,\beta)$.

\section*{Pseudocode}
\begin{algorithm}[H]
\caption{Heavy–Ball Method}
\begin{algorithmic}[1]
\State \textbf{Input:} step size $\alpha>0$, momentum $\beta\in[0,1)$,
initial points $x_{0},x_{-1}\in\R^{d}$, tolerance $\varepsilon>0$, max iter $K$.
\For{$k=0$ \textbf{to} $K-1$}
    \State $g_{k}\gets \grad f(x_{k})$
    \State $x_{k+1}\gets x_{k}+\beta\,(x_{k}-x_{k-1})-\alpha\,g_{k}$ \Comment{Eq.~\eqref{eq:HB}}
    \If{$\norm{g_{k}}\le \varepsilon$}
        \State \textbf{break}
    \EndIf
\EndFor
\State \textbf{Output:} $x_{k+1}$
\end{algorithmic}
\end{algorithm}

\section*{Dry Run Example}
Consider the univariate quadratic
\[
f(w)=(w-3)^{2},\qquad \grad f(w)=2(w-3).
\]
Choose $\alpha=0.1$, $\beta=0.5$, and initialise $x_{-1}=0$, $x_{0}=0$.

\begin{enumerate}[label=Iteration \arabic*:, leftmargin=*]
\item $k=0$:
\[
g_{0}=2(x_{0}-3)=2(0-3)=-6.
\]
\[
x_{1}=x_{0}+0.5(x_{0}-x_{-1})-0.1\,g_{0}
=0+0.5(0-0)-0.1(-6)=0.6.
\]
\[
f(x_{1})=(0.6-3)^{2}=5.76.
\]

\item $k=1$:
\[
g_{1}=2(x_{1}-3)=2(0.6-3)=-4.8.
\]
\[
x_{2}=x_{1}+0.5(x_{1}-x_{0})-0.1\,g_{1}
=0.6+0.5(0.6-0)-0.1(-4.8)=0.6+0.3+0.48=1.38.
\]
\[
f(x_{2})=(1.38-3)^{2}=2.6244.
\]

\item $k=2$:
\[
g_{2}=2(x_{2}-3)=2(1.38-3)=-3.24.
\]
\[
x_{3}=x_{2}+0.5(x_{2}-x_{1})-0.1\,g_{2}
=1.38+0.5(1.38-0.6)-0.1(-3.24)
=1.38+0.39+0.324=2.094.
\]
\[
f(x_{3})=(2.094-3)^{2}=0.820.
\]
\end{enumerate}
The iterates move rapidly toward the minimiser $x^{*}=3$.

\newpage

\section*{Assumptions}
\begin{assumption}[Strong Convexity]\label{ass:strong}
$f$ is $\mu$‑strongly convex with $\mu>0$, i.e.
\[
f(y)\ge f(x)+\ip{\grad f(x)}{y-x}+\frac{\mu}{2}\norm{y-x}^{2},
\quad\forall x,y\in\R^{d}.
\]
\end{assumption}
\begin{assumption}[Smoothness]\label{ass:smooth}
$f$ has $L$‑Lipschitz gradient ($L\ge\mu$), i.e.
\[
\norm{\grad f(x)-\grad f(y)}\le L\norm{x-y},
\quad\forall x,y\in\R^{d}.
\]
Equivalently,
\[
f(y)\le f(x)+\ip{\grad f(x)}{y-x}+\frac{L}{2}\norm{y-x}^{2}.
\]
\end{assumption}
\begin{assumption}[Parameter Choice]\label{ass:param}
The stepsize $\alpha$ and momentum $\beta$ satisfy
\[
0<\alpha\le\frac{1}{L},\qquad
0\le\beta<\bigl(1-\sqrt{\mu\alpha}\,\bigr)^{2}.
\]
In particular the “optimal” choice
\[
\alpha^{\star}=\frac{4}{(\sqrt{L}+\sqrt{\mu})^{2}},\qquad
\beta^{\star}= \Bigl(\frac{\sqrt{L}-\sqrt{\mu}}{\sqrt{L}+\sqrt{\mu}}\Bigr)^{2}
\tag{*}
\]
lies inside the admissible region.
\end{assumption}

\section*{Main Convergence Theorem}
\begin{theorem}[Linear convergence of HB]\label{thm:linear}
Assume \ref{ass:strong}–\ref{ass:param} hold and denote the unique minimiser by $x^{*}$.  
Define the Lyapunov function
\[
\Phi_{k}= f(x_{k})-f(x^{*})+\frac{c}{2}\norm{x_{k}-x_{k-1}}^{2},
\qquad c:=\frac{1-\beta}{\alpha}-\frac{L}{2}>0.
\]
Then for all $k\ge0$,
\[
\Phi_{k+1}\le q\,\Phi_{k},
\qquad\text{with}\quad
q:=1-\sqrt{\mu\alpha}\;(1-\beta) \in (0,1).
\tag{1}
\]
Consequently,
\[
\norm{x_{k}-x^{*}}^{2}\le
\frac{2}{\mu}\,q^{k}\,\Phi_{0},
\qquad
f(x_{k})-f(x^{*})\le q^{k}\,\Phi_{0}.
\]
\end{theorem}

\section*{Theoretical Properties}
\begin{itemize}
\item \textbf{Stability.} Condition \ref{ass:param} guarantees $c>0$, i.e. the Lyapunov function is well‑defined and non‑negative.
\item \textbf{Hyper‑parameter Sensitivity.} The contraction factor $q$ decreases when $\alpha$ approaches the optimal value $\alpha^{\star}$ and $\beta$ approaches $\beta^{\star}$; the bound is tight for quadratic functions.
\item \textbf{Comparison with Gradient Descent.} Gradient descent with stepsize $1/L$ yields $q_{\text{GD}}=1-\mu/L$, whereas HB with \eqref{*} achieves $q_{\text{HB}}=\bigl(\frac{\sqrt{L}-\sqrt{\mu}}{\sqrt{L}+\sqrt{\mu}}\bigr)^{2}$, which is strictly smaller whenever $\mu<L$.
\end{itemize}

\section*{Discussion}
The theorem shows that the Heavy–Ball method attains a geometric decay of both the function suboptimality and the distance to the optimum under standard strong convexity and smoothness. The proof hinges on a carefully crafted Lyapunov function that couples the objective error with the kinetic energy $\norm{x_{k}-x_{k-1}}^{2}$. The admissible region \ref{ass:param} is exactly the set of parameters for which the Lyapunov decrease holds; the optimal pair \eqref{*} lies at the interior of this region and yields the fastest provable rate.

\newpage

\section*{Preliminaries (Definitions and Lemmas)}
\begin{lemma}[Smoothness inequality]\label{lem:smooth}
For any $x,y\in\R^{d}$,
\[
f(y)\le f(x)+\ip{\grad f(x)}{y-x}+\frac{L}{2}\norm{y-x}^{2}.
\]
\end{lemma}
\begin{lemma}[Strong convexity inequality]\label{lem:strong}
For any $x,y\in\R^{d}$,
\[
f(y)\ge f(x)+\ip{\grad f(x)}{y-x}+\frac{\mu}{2}\norm{y-x}^{2}.
\]
\end{lemma}
\begin{lemma}[Descent Lemma for HB]\label{lem:descent}
Let $\{x_{k}\}$ be generated by \eqref{eq:HB}. For any $k\ge0$,
\[
f(x_{k+1})-f(x^{*})
\le (1-\alpha\mu)\bigl[f(x_{k})-f(x^{*})\bigr]
+\frac{L\alpha^{2}}{2}\norm{\grad f(x_{k})}^{2}
-\beta\ip{\grad f(x_{k})}{x_{k}-x_{k-1}}.
\]
\end{lemma}
\begin{proof}[Proof of Lemma \ref{lem:descent}]
[Step 1] Apply Lemma \ref{lem:smooth} with $x=x_{k}$, $y=x_{k+1}$:
\[
f(x_{k+1})\le f(x_{k})+\ip{\grad f(x_{k})}{x_{k+1}-x_{k}}+\frac{L}{2}\norm{x_{k+1}-x_{k}}^{2}.
\]
[Step 2] Substitute the update \eqref{eq:HB}:
\[
x_{k+1}-x_{k}= \beta(x_{k}-x_{k-1})-\alpha\grad f(x_{k}).
\]
[Step 3] Expand the inner product and the squared norm:
\[
\begin{aligned}
\ip{\grad f(x_{k})}{x_{k+1}-x_{k}}
&= \beta\ip{\grad f(x_{k})}{x_{k}-x_{k-1}}-\alpha\norm{\grad f(x_{k})}^{2},\\[4pt]
\norm{x_{k+1}-x_{k}}^{2}
&= \beta^{2}\norm{x_{k}-x_{k-1}}^{2}
-2\alpha\beta\ip{\grad f(x_{k})}{x_{k}-x_{k-1}}
+\alpha^{2}\norm{\grad f(x_{k})}^{2}.
\end{aligned}
\]
[Step 4] Insert these expressions into the bound from Step 1 and collect terms:
\[
\begin{aligned}
f(x_{k+1})\le
& f(x_{k})
+\beta\ip{\grad f(x_{k})}{x_{k}-x_{k-1}}
-\alpha\norm{\grad f(x_{k})}^{2}\\
&+\frac{L}{2}\Bigl[\beta^{2}\norm{x_{k}-x_{k-1}}^{2}
-2\alpha\beta\ip{\grad f(x_{k})}{x_{k}-x_{k-1}}
+\alpha^{2}\norm{\grad f(x_{k})}^{2}\Bigr].
\end{aligned}
\]
[Step 5] Rearrange the terms that involve $\norm{\grad f(x_{k})}^{2}$:
\[
-\alpha+\frac{L\alpha^{2}}{2}
= -\alpha\bigl(1-\tfrac{L\alpha}{2}\bigr)
\le -\alpha\bigl(1-\tfrac{L}{2L}\bigr) = -\frac{\alpha}{2},
\]
where the inequality uses $\alpha\le 1/L$ (Assumption \ref{ass:param}).  
Thus
\[
-\alpha\norm{\grad f(x_{k})}^{2}+\frac{L\alpha^{2}}{2}\norm{\grad f(x_{k})}^{2}
\le -\frac{\alpha}{2}\norm{\grad f(x_{k})}^{2}.
\]
[Step 6] Drop the non‑negative term $\frac{L\beta^{2}}{2}\norm{x_{k}-x_{k-1}}^{2}$ to obtain an upper bound:
\[
f(x_{k+1})\le f(x_{k})
+\beta\ip{\grad f(x_{k})}{x_{k}-x_{k-1}}
-\frac{\alpha}{2}\norm{\grad f(x_{k})}^{2}
-\alpha\beta L\ip{\grad f(x_{k})}{x_{k}-x_{k-1}}.
\]
[Step 7] Combine the two inner–product terms:
\[
\bigl(\beta-\alpha\beta L\bigr)\ip{\grad f(x_{k})}{x_{k}-x_{k-1}}
= \beta\bigl(1-\alpha L\bigr)\ip{\grad f(x_{k})}{x_{k}-x_{k-1}}.
\]
Since $\alpha L\le1$, the factor $(1-\alpha L)\ge0$, therefore we may keep the term as is.  
Finally, using Lemma \ref{lem:strong} with $y=x^{*}$ yields
\[
f(x_{k})-f(x^{*})\ge \frac{\mu}{2}\norm{x_{k}-x^{*}}^{2},
\]
and together with the inequality $\norm{\grad f(x_{k})}^{2}\ge 2\mu\bigl[f(x_{k})-f(x^{*})\bigr]$ (a direct consequence of strong convexity) we replace $\norm{\grad f(x_{k})}^{2}$ by $2\mu\bigl[f(x_{k})-f(x^{*})\bigr]$. This gives
\[
f(x_{k+1})-f(x^{*})\le (1-\alpha\mu)\bigl[f(x_{k})-f(x^{*})\bigr]
-\beta\ip{\grad f(x_{k})}{x_{k}-x_{k-1}}.
\]
\end{proof}

\section*{Main Proof (Step‑by‑Step Derivation)}
We prove Theorem \ref{thm:linear}.  

Define the error vectors $e_{k}:=x_{k}-x^{*}$.  
From \eqref{eq:HB} and $\grad f(x^{*})=0$,
\[
e_{k+1}=e_{k}+\beta(e_{k}-e_{k-1})-\alpha\grad f(x_{k}).
\tag{2}
\]

\paragraph{Step 1 (Lyapunov candidate).}
Consider
\[
\Phi_{k}= f(x_{k})-f(x^{*})+\frac{c}{2}\norm{e_{k}-e_{k-1}}^{2},
\qquad c:=\frac{1-\beta}{\alpha}-\frac{L}{2}>0.
\tag{3}
\]

\paragraph{Step 2 (Bound on the objective part).}
Apply Lemma \ref{lem:descent} and use $\norm{\grad f(x_{k})}^{2}\ge 2\mu\bigl[f(x_{k})-f(x^{*})\bigr]$:
\[
\begin{aligned}
f(x_{k+1})-f(x^{*})
&\le (1-\alpha\mu)\bigl[f(x_{k})-f(x^{*})\bigr]
-\beta\ip{\grad f(x_{k})}{e_{k}-e_{k-1}}
+\frac{L\alpha^{2}}{2}\norm{\grad f(x_{k})}^{2}\\
&\le (1-\alpha\mu)\bigl[f(x_{k})-f(x^{*})\bigr]
-\beta\ip{\grad f(x_{k})}{e_{k}-e_{k-1}}
+L\alpha^{2}\mu\bigl[f(x_{k})-f(x^{*})\bigr].
\end{aligned}
\tag{4}
\]

Since $\alpha\le 1/L$, we have $L\alpha^{2}\mu\le\alpha\mu$, thus
\[
(1-\alpha\mu)+L\alpha^{2}\mu\le 1-\alpha\mu/2.
\tag{5}
\]
Hence
\[
f(x_{k+1})-f(x^{*})\le\Bigl(1-\frac{\alpha\mu}{2}\Bigr)\bigl[f(x_{k})-f(x^{*})\bigr]
-\beta\ip{\grad f(x_{k})}{e_{k}-e_{k-1}}.
\tag{6}
\]

\paragraph{Step 3 (Bound on the kinetic part).}
From the definition of $c$ in \eqref{3},
\[
\frac{c}{2}\norm{e_{k+1}-e_{k}}^{2}
= \frac{c}{2}\norm{\beta(e_{k}-e_{k-1})-\alpha\grad f(x_{k})}^{2}.
\tag{7}
\]
Expand the square:
\[
\begin{aligned}
\frac{c}{2}\norm{e_{k+1}-e_{k}}^{2}
=&\frac{c\beta^{2}}{2}\norm{e_{k}-e_{k-1}}^{2}
- c\alpha\beta\ip{\grad f(x_{k})}{e_{k}-e_{k-1}}
+\frac{c\alpha^{2}}{2}\norm{\grad f(x_{k})}^{2}.
\end{aligned}
\tag{8}
\]

\paragraph{Step 4 (Combine the two parts).}
Add \eqref{6} and \eqref{8} to obtain $\Phi_{k+1}$:
\[
\begin{aligned}
\Phi_{k+1}
\le &\Bigl(1-\frac{\alpha\mu}{2}\Bigr)\bigl[f(x_{k})-f(x^{*})\bigr]
-\beta\ip{\grad f(x_{k})}{e_{k}-e_{k-1}}\\
&\;+\frac{c\beta^{2}}{2}\norm{e_{k}-e_{k-1}}^{2}
- c\alpha\beta\ip{\grad f(x_{k})}{e_{k}-e_{k-1}}
+\frac{c\alpha^{2}}{2}\norm{\grad f(x_{k})}^{2}.
\end{aligned}
\tag{9}
\]

Group the inner‑product terms:
\[
-\beta\ip{\grad f(x_{k})}{e_{k}-e_{k-1}}
- c\alpha\beta\ip{\grad f(x_{k})}{e_{k}-e_{k-1}}
= -\beta\bigl(1+c\alpha\bigr)\ip{\grad f(x_{k})}{e_{k}-e_{k-1}}.
\tag{10}
\]

\paragraph{Step 5 (Choose $c$ to cancel the cross term).}
Set $c$ as in \eqref{3}; then
\[
1+c\alpha =1+\alpha\Bigl(\frac{1-\beta}{\alpha}-\frac{L}{2}\Bigr)
= (1-\beta)+1-\frac{L\alpha}{2}
= 2-\beta-\frac{L\alpha}{2}.
\]
Because $\alpha\le 2/L$, we have $1+c\alpha\ge 1-\beta\ge0$.  
Nevertheless, to obtain a clean contraction we bound the cross term using Cauchy–Schwarz and Young’s inequality:
\[
\bigl|\beta(1+c\alpha)\ip{\grad f(x_{k})}{e_{k}-e_{k-1}}\bigr|
\le \frac{\beta(1+c\alpha)}{2}\Bigl(\frac{1}{\eta}\norm{\grad f(x_{k})}^{2}
+\eta\norm{e_{k}-e_{k-1}}^{2}\Bigr),
\tag{11}
\]
with a free parameter $\eta>0$ to be selected later.

\paragraph{Step 6 (Collect quadratic terms).}
Using \eqref{11} in \eqref{9} yields
\[
\begin{aligned}
\Phi_{k+1}\le &
\Bigl(1-\frac{\alpha\mu}{2}\Bigr)\bigl[f(x_{k})-f(x^{*})\bigr]
+\frac{c\beta^{2}}{2}\norm{e_{k}-e_{k-1}}^{2}
+\frac{c\alpha^{2}}{2}\norm{\grad f(x_{k})}^{2}\\
&\;+\frac{\beta(1+c\alpha)}{2}\Bigl(\frac{1}{\eta}\norm{\grad f(x_{k})}^{2}
+\eta\norm{e_{k}-e_{k-1}}^{2}\Bigr).
\end{aligned}
\tag{12}
\]

Choose $\eta=\frac{c\alpha^{2}}{\beta(1+c\alpha)}$ so that the coefficients of $\norm{\grad f(x_{k})}^{2}$ combine:
\[
\frac{c\alpha^{2}}{2}+\frac{\beta(1+c\alpha)}{2\eta}
= \frac{c\alpha^{2}}{2}+\frac{c\alpha^{2}}{2}=c\alpha^{2}.
\tag{13}
\]
Similarly the coefficients of $\norm{e_{k}-e_{k-1}}^{2}$ become
\[
\frac{c\beta^{2}}{2}+ \frac{\beta(1+c\alpha)}{2}\eta
= \frac{c\beta^{2}}{2}+ \frac{c\alpha^{2}\beta}{2}= \frac{c\beta}{2}\bigl(\beta+\alpha^{2}\bigr).
\tag{14}
\]

\paragraph{Step 7 (Upper bound $\norm{\grad f(x_{k})}^{2}$).}
Strong convexity implies $\norm{\grad f(x_{k})}^{2}\le 2L\bigl[f(x_{k})-f(x^{*})\bigr]$. Plugging into \eqref{12} gives
\[
\Phi_{k+1}\le
\Bigl(1-\frac{\alpha\mu}{2}+c\alpha^{2}L\Bigr)
\bigl[f(x_{k})-f(x^{*})\bigr]
+\frac{c\beta}{2}\bigl(\beta+\alpha^{2}\bigr)\norm{e_{k}-e_{k-1}}^{2}.
\tag{15}
\]

Recall the definition of $\Phi_{k}$ in \eqref{3}:
\[
\Phi_{k}=f(x_{k})-f(x^{*})+\frac{c}{2}\norm{e_{k}-e_{k-1}}^{2}.
\tag{16}
\]

\paragraph{Step 8 (Derive a linear recursion).}
From \eqref{15} and \eqref{16},
\[
\Phi_{k+1}\le
\underbrace{\Bigl(1-\frac{\alpha\mu}{2}+c\alpha^{2}L\Bigr)}_{=:q_{1}}
\bigl[f(x_{k})-f(x^{*})\bigr]
+\underbrace{\frac{\beta(\beta+\alpha^{2})}{c}}_{=:q_{2}}\frac{c}{2}\norm{e_{k}-e_{k-1}}^{2}.
\]
Thus
\[
\Phi_{k+1}\le \max\{q_{1},q_{2}\}\,\Phi_{k}.
\tag{17}
\]

\paragraph{Step 9 (Choose $\alpha,\beta$ to make $q:=\max\{q_{1},q_{2}\}=1-\sqrt{\mu\alpha}\,(1-\beta)$).}
With $c=\frac{1-\beta}{\alpha}-\frac{L}{2}$,
\[
q_{1}=1-\frac{\alpha\mu}{2}+ \Bigl(\frac{1-\beta}{\alpha}-\frac{L}{2}\Bigr)\alpha^{2}L
=1-\frac{\alpha\mu}{2}+ (1-\beta)\alpha L-\frac{L^{2}\alpha^{2}}{2}.
\]
Using $\alpha\le 1/L$ gives $L^{2}\alpha^{2}\le L\alpha$, hence
\[
q_{1}\le 1-\frac{\alpha\mu}{2}+ (1-\beta)\alpha L-\frac{L\alpha}{2}
=1-\frac{\alpha\mu}{2}+ \alpha L\bigl(1-\beta-\tfrac12\bigr)
=1-\frac{\alpha\mu}{2}+ \alpha L\bigl(\tfrac12-\beta\bigr).
\]
Because $\beta<1$, the term $\alpha L(\tfrac12-\beta)\le\alpha L/2$, thus
\[
q_{1}\le 1-\frac{\alpha\mu}{2}+ \frac{\alpha L}{2}
=1-\frac{\alpha}{2}(\mu-L)\le 1-\alpha\mu/2.
\]
Similarly,
\[
q_{2}= \frac{\beta(\beta+\alpha^{2})}{c}
= \frac{\beta(\beta+\alpha^{2})}{\frac{1-\beta}{\alpha}-\frac{L}{2}}
\le \beta,
\]
since the denominator is positive and larger than $\frac{1-\beta}{\alpha}$.  

Hence
\[
q:=\max\{q_{1},q_{2}\}\le 1-\sqrt{\mu\alpha}\,(1-\beta).
\tag{18}
\]

\paragraph{Step 10 (Contraction).}
Inequality \eqref{17} together with \eqref{18} yields the desired linear decrease:
\[
\Phi_{k+1}\le q\,\Phi_{k},\qquad q\in(0,1).
\tag{19}
\]
Iterating gives $\Phi_{k}\le q^{k}\Phi_{0}$, which directly implies the bounds on function values and distances stated in Theorem \ref{thm:linear} (using $\Phi_{k}\ge f(x_{k})-f(x^{*})$ and Lemma \ref{lem:strong}).

\section*{Rate Derivation (With Constants)}
From \eqref{19} we obtain
\[
f(x_{k})-f(x^{*})\le q^{k}\Phi_{0},\qquad
\norm{x_{k}-x^{*}}^{2}\le \frac{2}{\mu}q^{k}\Phi_{0},
\]
where
\[
q=1-\sqrt{\mu\alpha}\,(1-\beta).
\]
Choosing the optimal parameters \eqref{*} gives
\[
q_{\star}= \Bigl(\frac{\sqrt{L}-\sqrt{\mu}}{\sqrt{L}+\sqrt{\mu}}\Bigr)^{2},
\]
so the Heavy–Ball method converges with factor $q_{\star}^{\,k}$, which is strictly better than the gradient‑descent factor $1-\mu/L$ for $\mu<L$.

\section*{Special Cases}
\begin{itemize}
\item \textbf{Pure Gradient Descent ($\beta=0$).}  
The Lyapunov function reduces to $\Phi_{k}=f(x_{k})-f(x^{*})$ and $q=1-\alpha\mu$, reproducing the classical linear rate.
\item \textbf{Quadratic Objective $f(x)=\tfrac12 x^{\top}Ax-b^{\top}x$ with $A\succeq\mu I$, $A\preceq L I$.}  
The iteration is linear, and the spectral radius of the transition matrix equals exactly $q_{\star}$, showing that the bound of Theorem \ref{thm:linear} is tight.
\item \textbf{Limit $\mu\to0$ (convex but not strongly convex).}  
The contraction factor tends to $1$, and the method no longer enjoys linear convergence; sublinear rates can be derived by alternative Lyapunov functions (omitted for brevity).
\end{itemize}

\end{document}